\SourceLeaf{72}%
a concise overview of all the remaining operations needed for a first determination of a comet’s orbit. We note that some of the other formulas also appear here in a form different from the usual one. The calculation is a little shorter in that form only if one uses the table the author published a year ago for calculating the logarithms of sums and differences directly. In this way everything needed to calculate a comet’s orbit, apart from the tables, is gathered here on two or three pages. The whole procedure is also illustrated by an example using the latest comet, in which the author derives elements from his observations of 7, 14 and 21 April. They differ only slightly from the elements given above, which are based on all the observations.
\par\bigskip\centerline{\rule{.24\textwidth}{.3pt}}
